\documentclass[11pt]{article}
\usepackage{mathblatt}
% Präambel des Setzers (werkzeuge/setzer-praeambel.tex, wortgleich); dazu \gteil (Muster T2U)
\makeatletter
\setlength{\mbpfbe}{0mm}% keine Punkte (pruefung.md § 5: nur beim Prüfstein)
\newlength{\gEin}\setlength{\gEin}{0pt}
\let\gSchrift\relax
\renewcommand{\pfkreuzzeile}[1]{\par\noindent\hangindent3.2em\hangafter1\hspace*{1.6em}$\square$\ #1\par}
\newcommand{\gkreuz}[1]{$\square$\ #1\hspace{2em}}
\newcommand{\gkreuzzeile}[2]{\par\noindent\hangindent2.2em\hangafter1\mbox{#1\,$\square$}\ #2\par}% Buchstabe, Option: eine je Zeile (neben Skizze)
% links, was man ansieht, rechts, was man tut; Spaltengrenze überall an derselben Stelle (gemeinsam.md „Satz“)
\newsavebox{\gbox}\newlength{\gLinks}
% v0.10 (Lehrer 13.10.): keine starre Halbierung – die linke Spalte ist so breit wie die Figur (+ 6 mm), die
% Antwort steht direkt daneben; ist die Figur breiter als 0,62 Zeile, steht die Antwort darunter
\newcommand{\gzwei}[2]{\par\smallskip\noindent\sbox{\gbox}{#1}%
  \setlength{\gLinks}{\dimexpr\wd\gbox+6mm\relax}%
  \ifdim\gLinks>0.62\linewidth
    \usebox{\gbox}\par\smallskip\noindent\begin{minipage}{\linewidth}#2\end{minipage}%
  \else
    \begin{minipage}[t]{\gLinks}\vspace{0pt}\usebox{\gbox}\end{minipage}%
    \begin{minipage}[t]{\dimexpr\linewidth-\gLinks\relax}\vspace{0pt}\raggedright #2\end{minipage}%
  \fi\par}
% v0.9 (12.10., pruefung.md § 5 „Marke mit B, Bedingungen rechts“): drittes Argument von pfaufg = Bedingungen
% der Prüfung (v0.10: \pfohnetr = TR durchgestrichen, „FOR“), klein grau im rechten Rand auf der Grundlinie der Nummer, Zeilen mit \\
\newcommand{\gbedrand}[2]{\ifblank{#2}{}{\rlap{\hspace*{\dimexpr#1+2mm\relax}%
  \parbox[t]{15mm}{\raggedright\scriptsize\leavevmode\color{mbgrau}\strut#2}}}}
% Teil eines Bündels (a, b …): Marke links im Rand, Bedingungen rechts im Rand, auf der Zeile des Teils
\newcommand{\gteilrand}[2]{\llap{\mbpfrandmarke{#1}\hspace*{\mbpfnr}}\gbedrand{\linewidth}{#2}}
% Teil eines Bündels: Buchstabe in eigener Spalte, Inhalt eingerückt (\gEin = Einzug, damit die Spaltengrenze von
% \gzwei an derselben Stelle bleibt)
\newenvironment{gbteil}[3]{\par\smallskip\noindent\gteilrand{#2}{#3}\makebox[\mbpfteil][l]{\textbf{#1}}%
  \begin{minipage}[t]{\dimexpr\linewidth-\mbpfteil\relax}\setlength{\gEin}{\mbpfteil}\raggedright\strut\ignorespaces}%
 {\end{minipage}\par}
% Grafik höchstens so breit wie die Spalte (Bündel nebeneinander, Koordinatensystem ohne Karo)
\newsavebox{\gpbox}
\newcommand{\gpassend}[1]{\sbox{\gpbox}{#1}\ifdim\wd\gpbox>\linewidth\typeout{GPASSEND: verkleinert \the\wd\gpbox\space auf \the\linewidth}\resizebox{\linewidth}{!}{\usebox{\gpbox}}\else\usebox{\gpbox}\fi}
% Skizzenfeld (v0.10, Lehrer 13.10.): Karo ab dem linken Rand über die Zeile, „Skizze“ klein grau oben links im Feld
% (wie \pfskizzenfeld der Vorlage); der rechte Rand bleibt den Bedingungen
\newcommand{\gskizze}[1]{\par\smallskip\noindent\begin{tikzpicture}
  \clip (0,0) rectangle ({\linewidth-1mm},{#1*10mm});
  \draw[black!16,line width=0.3pt,step=5mm] (0,0) grid ({\linewidth},{#1*10mm});
  \node[anchor=north west,inner sep=2pt,font=\scriptsize,text=mbgrau,fill=white] at (0,{#1*10mm}) {Skizze};
  \end{tikzpicture}\par}
\RenewDocumentEnvironment{pfaufg}{m m m}{%
  \begin{lrbox}{\mbaufgabebox}\begin{minipage}[t]{\linewidth}%
  \gSchrift\noindent\hspace*{\gEin}\mbpfrandmarke{#1}\gbedrand{\linewidth-\gEin}{#3}%
  \begin{minipage}[t]{\mbpfnr}\strut\textbf{#2}\end{minipage}%
  \begin{minipage}[t]{\dimexpr\linewidth-\gEin-\mbpfnr\relax}\raggedright\strut\ignorespaces}%
 {\end{minipage}%
  \end{minipage}\end{lrbox}%
  \setlength{\mbaufgabehoehe}{\dimexpr\ht\mbaufgabebox+\dp\mbaufgabebox\relax}%
  \ifdefempty{\gkopfnach}{\par\addvspace{5pt}\Needspace*{\mbaufgabehoehe}}{%
    % Stufenkopf, der mit dieser (ersten) Aufgabe zusammenbleibt: Platz für beide verlangen, dann den Kopf setzen
    % (12.10., Quadratische: ein Kopf stand allein unten, die Aufgabe mit Koordinatensystem auf der nächsten Seite)
    \par\Needspace*{\dimexpr\mbaufgabehoehe+40pt\relax}\gkopfnach\gdef\gkopfnach{}\par\addvspace{5pt}}%
  \noindent\usebox{\mbaufgabebox}\par\addvspace{5pt}}
% Stufenkopf vormerken: \gkopfvor{<Kopf>} – die nächste pfaufg setzt ihn, sobald beide zusammen Platz haben
\def\gkopfnach{}
\newcommand{\gkopfvor}[1]{\gdef\gkopfnach{#1}}
% Teilaufgabe mit Skizze links in gzwei: Buchstabe in eigener Spalte, darunter Text und Skizze
\newcommand{\gteil}[2]{\makebox[\mbpfteil][l]{\textbf{#1}}\parbox[t]{\dimexpr0.5\textwidth-\gEin-\mbpfnr-5mm-\mbpfteil\relax}{\raggedright #2}}
% Stufenkopf der Vorlage (mathblatt.sty 2026-10-12: Linie, Abstand, 13 pt, \pfjahre) ohne ihren Vorhalt von
% 14 Zeilen: der Setzer fordert vor dem Kopf 0,16 Seitenhöhe (Kopf + erste Aufgabe), mehr nicht
\setlength{\mbpfstufevorhalt}{0pt}
\makeatother
\newcommand{\karo}[1]{\par\smallskip\noindent\begin{tikzpicture}
  \clip (0,0) rectangle ({\linewidth-1mm},{#1*10mm});
  \draw[black!16,line width=0.3pt,step=5mm] (0,0) grid ({\linewidth},{#1*10mm});
  \end{tikzpicture}\par}
\newcommand{\kk}[1]{$\square$\,#1\hspace{0.9em}}
\newcommand{\linie}[1]{\,{\color{black!45}\rule[-1pt]{#1}{0.4pt}}\,}
\newcommand{\pkt}{\textperiodcentered{}}
% v0.14 (14.10.): Stufenkopf mit fünf Punkten (Option = Jahre für die Bestellkarte) oder „zuletzt 19“
\NewDocumentCommand{\gpunkte}{o m}{#2}
\newcommand{\gzuletzt}[1]{zuletzt #1}
% v0.7 (12.10.): Schülerfrage unter dem Stufenkopf; Ankreuzoption als Bild (Buchstabe und Kästchen darunter)
\newcommand{\gfrage}[1]{\par\nobreak\noindent\hspace*{\mbpfnr}{\small\itshape #1}\par\nobreak\addvspace{2pt}}
\newlength{\gbildbreite}\setlength{\gbildbreite}{0.32\linewidth}% Setzer setzt sie je Reihe (vier Bilder: 0,24)
\newcommand{\gbildoption}[3]{\begin{minipage}[t][#3][b]{\gbildbreite}\centering #2\par\smallskip\mbox{#1\,$\square$}\end{minipage}}
% Lineare (12.10.): Bild mit Buchstabe ohne Kästchen (Diagramm zuordnen); Aussage mit „wahr / falsch“ zum Ankreuzen
\newcommand{\gbildname}[3]{\begin{minipage}[t][#3][b]{\gbildbreite}\centering #2\par\smallskip\mbox{\textbf{#1}}\end{minipage}}
\newcommand{\gwahrfalsch}[2]{\par\noindent\hangindent2.2em\hangafter1\makebox[2.2em][l]{(#1)}#2\hfill\mbox{wahr $\square$\quad falsch $\square$}\par}
\begin{document}

\pfheftstil
\pflaufkopf[Abitur GK \pkt{} Prüfung]{Säulendiagramm und Erwartungswert}{AS80}\pfseitevon
\pfheftkopf{Säulendiagramm und Erwartungswert}{}
\pfstufenzeile{Wahrscheinlichkeit ablesen \pkt{} Erwartungswert am Diagramm \pkt{} Erwartungswert berechnen}

% ===================== Stufe: Wahrscheinlichkeit ablesen =====================
\gkopfvor{\pfstufekopf{Wahrscheinlichkeit ablesen}{\gpunkte[26 25 23]{\textcolor{black!80}{$\bullet$}\,\textcolor{black!80}{$\bullet$}\,\textcolor{black!25}{$\circ$}\,\textcolor{black!80}{$\bullet$}\,\textcolor{black!25}{$\circ$}}}\special{pdf:dest (stufe-wahrscheinlichkeit-ablesen) [@thispage /XYZ @xpos @ypos null]}}
% kein Einstieg: keine Schülerfrage in der Gliederung (Schülerfrage, -Fälle, -Form, -Weg, -Skizze; oder Frage) – Einstieg nur von dort, nie aus der Bank (pruefung.md § 5, 12.10.)
% 2023-bebb-gk-A1.7a-m1 (Original 2023-bebb-gk-A1.7a, Schwere I, rang 1; herausgelöste Fassung: keine)
\begin{pfaufg}{GK~’23\\A1.7a}{1.}{\pfohnetr}
Die Abbildung zeigt die Wahrscheinlichkeitsverteilung einer binomialverteilten Zufallsgröße $X$. Bestimme mit der Abbildung einen Näherungswert für $P(4 < X < 7)$.
\gzwei{\saeulenab[ymax=0.3,ystep=0.05,yfein=0.01,ylabel={$P(X=k)$},hoehe=2.8,breite=0.55]{0/0.0060,1/0.0403,2/0.1209,3/0.2150,4/0.2508,5/0.2007,6/0.1115,7/0.0425,8/0.0106,9/0.0016,10/0.0001}}{\karo{2}}
\fusshilfe{\mbox{1:~0,31}}
\end{pfaufg}
% 2026-bb-gk-A1.6b-m1 (Original 2026-bb-gk-A1.6b, Schwere II, rang 2; herausgelöste Fassung: keine)
\begin{pfaufg}{GK~’26\\A1.6b}{2.}{\pfohnetr}
Die Abbildung zeigt die Wahrscheinlichkeitsverteilung einer binomialverteilten Zufallsgröße $X$ mit $n = 16$ und unbekannter Trefferwahrscheinlichkeit $p$. Gib mit der Abbildung die natürlichen Zahlen $v$ und $w$ an, für die gilt: $P(X = v) \approx 0{,}18$ und $0{,}15 < P(X \le w) < 0{,}25$.
\gzwei{\saeulenab[ymax=0.25,ystep=0.05,yfein=0.01,ylabel={$P(X=k)$},hoehe=2.8,breite=0.5]{0/0.0100,1/0.0535,2/0.1336,3/0.2079,4/0.2252,5/0.1802,6/0.1101,7/0.0524,8/0.0197,9/0.0058,10/0.0014,11/0.0002,12/0.0000}}{\karo{3}}
\fusshilfe{\mbox{2:~5; 2}}
\end{pfaufg}
% 2026-bb-gk-B4c-m1 (Original 2026-bb-gk-B4c, Schwere II, rang 3; herausgelöste Fassung: keine)
\begin{pfaufg}{GK~’26\\B4c}{3.}{}
In einem Supermarkt sammeln 70 \% der Kundschaft Treuepunkte. Es werden 12 Personen aus der Kundschaft zufällig ausgewählt; die Zufallsgröße $X$ gibt an, wie viele davon Treuepunkte sammeln. $X$ wird als binomialverteilt angenommen. Die Abbildung zeigt die Wahrscheinlichkeitsverteilung der Zufallsgröße $Y$ mit $n = 12$ und $p = 0{,}3$. Erläutere, wie man mit der Abbildung die Wahrscheinlichkeit $P(X = 8)$ ermitteln kann.
\gzwei{\saeulenab[ymax=0.3,ystep=0.1,yfein=0.05,ylabel={$P(X=k)$},hoehe=2.8,breite=0.5]{0/0.0138,1/0.0712,2/0.1678,3/0.2397,4/0.2311,5/0.1585,6/0.0792,7/0.0291,8/0.0078,9/0.0015,10/0.0002,11/0.0000,12/0.0000}}{\karo{3}}
\fusshilfe{\mbox{3:~Säule von Y bei vier ablesen; 0,23}}
\end{pfaufg}
% 2025-bebb-gk-A1.9b-m1 (Original 2025-bebb-gk-A1.9b, Schwere III, rang 4; herausgelöste Fassung: keine)
\begin{pfaufg}{GK~’25\\A1.9b}{4.}{\pfohnetr}
Die Abbildung zeigt die Wahrscheinlichkeitsverteilung einer binomialverteilten Zufallsgröße $X$ mit der Trefferwahrscheinlichkeit $p = 0{,}5$ und unbekanntem $n$. Die beiden höchsten Säulen bei $k = 9$ und $k = 10$ sind gleich hoch. Es gilt $P(X \ge 8) \approx 0{,}82$ und $P(X = 11) \approx 0{,}14$. Berechne mit diesen Werten einen Näherungswert für $P(X = 9)$.
\gzwei{\saeulenab[ymax=0.2,ystep=0.05,yfein=0.05,ylabel={$P(X=k)$},hoehe=2.8,breite=0.5,ohnezahlen]{2/0.0003,3/0.0018,4/0.0074,5/0.0222,6/0.0518,7/0.0961,8/0.1442,9/0.1762,10/0.1762,11/0.1442,12/0.0961,13/0.0518,14/0.0222,15/0.0074,16/0.0018,17/0.0003}}{\karo{3}}
\fusshilfe{\mbox{4:~0,18}}
\end{pfaufg}

% ===================== Stufe: Erwartungswert am Diagramm =====================
\gkopfvor{\pfstufekopf{Erwartungswert am Diagramm}{\gpunkte[26 25 23]{\textcolor{black!80}{$\bullet$}\,\textcolor{black!80}{$\bullet$}\,\textcolor{black!25}{$\circ$}\,\textcolor{black!80}{$\bullet$}\,\textcolor{black!25}{$\circ$}}}\special{pdf:dest (stufe-erwartungswert-am-diagramm) [@thispage /XYZ @xpos @ypos null]}}
% kein Einstieg: keine Schülerfrage in der Gliederung (Schülerfrage, -Fälle, -Form, -Weg, -Skizze; oder Frage) – Einstieg nur von dort, nie aus der Bank (pruefung.md § 5, 12.10.)
% 2026-bb-gk-A1.6a-m1 (Original 2026-bb-gk-A1.6a, Schwere I, rang 5; herausgelöste Fassung: keine)
\begin{pfaufg}{GK~’26\\A1.6a}{5.}{\pfohnetr}
Die Abbildung zeigt die Wahrscheinlichkeitsverteilung einer binomialverteilten Zufallsgröße $X$ mit $n = 16$ und unbekannter Trefferwahrscheinlichkeit $p$. Der Erwartungswert von $X$ ist ganzzahlig. Ermittle mit der Abbildung den Wert von $p$.
\gzwei{\saeulenab[ymax=0.25,ystep=0.05,yfein=0.01,ylabel={$P(X=k)$},hoehe=2.8,breite=0.5]{0/0.0100,1/0.0535,2/0.1336,3/0.2079,4/0.2252,5/0.1802,6/0.1101,7/0.0524,8/0.0197,9/0.0058,10/0.0014,11/0.0002,12/0.0000}}{\karo{3}}
\fusshilfe{\mbox{5:~0,25}}
\end{pfaufg}
% 2023-bebb-gk-A1.7b-m1 (Original 2023-bebb-gk-A1.7b, Schwere II, rang 6; herausgelöste Fassung: keine)
\begin{pfaufg}{GK~’23\\A1.7b}{6.}{\pfohnetr}
Die Zufallsgröße $Y$ ist binomialverteilt mit $n = 15$ und $p = 0{,}2$. Genau eine der Abbildungen A, B und C zeigt die Wahrscheinlichkeitsverteilung von $Y$. Entscheide begründet, welche.
\par\smallskip\noindent\setlength{\gbildbreite}{\dimexpr0.96\linewidth/3\relax}\sbox{\gbox}{\vbox{\resizebox{0.92\gbildbreite}{!}{\raisebox{\depth}{\saeulenab[ymax=0.3,ystep=0.1,yfein=0.05,ylabel={$P(X=k)$},hoehe=2.8,breite=0.5]{0/0.0023,1/0.0171,2/0.0599,3/0.1299,4/0.1948,5/0.2143,6/0.1786,7/0.1148,8/0.0574,9/0.0223,10/0.0067,11/0.0015,12/0.0003,13/0.0000,14/0.0000,15/0.0000}}}\par\resizebox{0.92\gbildbreite}{!}{\raisebox{\depth}{\saeulenab[ymax=0.3,ystep=0.1,yfein=0.05,ylabel={$P(X=k)$},hoehe=2.8,breite=0.5]{0/0.0352,1/0.1319,2/0.2309,3/0.2501,4/0.1876,5/0.1032,6/0.0430,7/0.0138,8/0.0035,9/0.0007,10/0.0001,11/0.0000,12/0.0000,13/0.0000,14/0.0000,15/0.0000}}}\par\resizebox{0.92\gbildbreite}{!}{\raisebox{\depth}{\saeulenab[ymax=0.6,ystep=0.1,yfein=0.05,ylabel={$P(X=k)$},hoehe=2.8,breite=0.5]{0/0.0704,1/0.2639,2/0.4618,3/0.5003,4/0.3752,5/0.2064,6/0.0860,7/0.0276,8/0.0069,9/0.0013,10/0.0002,11/0.0000,12/0.0000,13/0.0000,14/0.0000,15/0.0000}}}}}\gbildoption{A}{\resizebox{0.92\gbildbreite}{!}{\raisebox{\depth}{\saeulenab[ymax=0.3,ystep=0.1,yfein=0.05,ylabel={$P(X=k)$},hoehe=2.8,breite=0.5]{0/0.0023,1/0.0171,2/0.0599,3/0.1299,4/0.1948,5/0.2143,6/0.1786,7/0.1148,8/0.0574,9/0.0223,10/0.0067,11/0.0015,12/0.0003,13/0.0000,14/0.0000,15/0.0000}}}}{\dimexpr\ht\gbox/3+\baselineskip\relax}\hfill\gbildoption{B}{\resizebox{0.92\gbildbreite}{!}{\raisebox{\depth}{\saeulenab[ymax=0.3,ystep=0.1,yfein=0.05,ylabel={$P(X=k)$},hoehe=2.8,breite=0.5]{0/0.0352,1/0.1319,2/0.2309,3/0.2501,4/0.1876,5/0.1032,6/0.0430,7/0.0138,8/0.0035,9/0.0007,10/0.0001,11/0.0000,12/0.0000,13/0.0000,14/0.0000,15/0.0000}}}}{\dimexpr\ht\gbox/3+\baselineskip\relax}\hfill\gbildoption{C}{\resizebox{0.92\gbildbreite}{!}{\raisebox{\depth}{\saeulenab[ymax=0.6,ystep=0.1,yfein=0.05,ylabel={$P(X=k)$},hoehe=2.8,breite=0.5]{0/0.0704,1/0.2639,2/0.4618,3/0.5003,4/0.3752,5/0.2064,6/0.0860,7/0.0276,8/0.0069,9/0.0013,10/0.0002,11/0.0000,12/0.0000,13/0.0000,14/0.0000,15/0.0000}}}}{\dimexpr\ht\gbox/3+\baselineskip\relax}\par
\karo{2}
\fusshilfe{\mbox{6:~B; B hat die höchste Säule bei 3, A bei 5, und bei C ist die Summe der Säulen größer als 1}}
\end{pfaufg}
% 2025-bebb-gk-A1.9a-m1 (Original 2025-bebb-gk-A1.9a, Schwere II, rang 7; herausgelöste Fassung: keine)
\begin{pfaufg}{GK~’25\\A1.9a}{7.}{\pfohnetr}
Die Abbildung zeigt die Wahrscheinlichkeitsverteilung einer binomialverteilten Zufallsgröße $X$ mit der Trefferwahrscheinlichkeit $p = 0{,}5$ und unbekanntem $n$. Die beiden höchsten Säulen bei $k = 9$ und $k = 10$ sind gleich hoch. Begründe, dass $n$ keine gerade Zahl ist.
\gzwei{\saeulenab[ymax=0.2,ystep=0.05,yfein=0.05,ylabel={$P(X=k)$},hoehe=2.8,breite=0.5,ohnezahlen]{2/0.0003,3/0.0018,4/0.0074,5/0.0222,6/0.0518,7/0.0961,8/0.1442,9/0.1762,10/0.1762,11/0.1442,12/0.0961,13/0.0518,14/0.0222,15/0.0074,16/0.0018,17/0.0003}}{\karo{3}}
\fusshilfe{\mbox{7:~19 ungerade}}
\end{pfaufg}
% 2025-bebb-gk-B4c-m1 (Original 2025-bebb-gk-B4c, Schwere II, rang 8; herausgelöste Fassung: keine)
\begin{pfaufg}{GK~’25\\B4c}{8.}{}
In einer Großstadt gelten 8,55 \% aller Haushalte als überbelegt. Es werden 800 Haushalte zufällig ausgewählt; die Zufallsgröße $X$ gibt die Anzahl der überbelegten Haushalte unter ihnen an und wird als binomialverteilt angenommen. Beurteile ohne Berechnung von Wahrscheinlichkeiten die Aussage: „Die Wahrscheinlichkeitsverteilung von $X$ nimmt für $k = 70$ ihren größten Wert an.“
\karo{2}
\fusshilfe{\mbox{8:~falsch}}
\end{pfaufg}

% ===================== Stufe: Erwartungswert berechnen =====================
\gkopfvor{\pfstufekopf{Erwartungswert berechnen}{\gpunkte[26 22]{\textcolor{black!80}{$\bullet$}\,\textcolor{black!25}{$\circ$}\,\textcolor{black!25}{$\circ$}\,\textcolor{black!25}{$\circ$}\,\textcolor{black!80}{$\bullet$}}}\special{pdf:dest (stufe-erwartungswert-berechnen) [@thispage /XYZ @xpos @ypos null]}}
% kein Einstieg: keine Schülerfrage in der Gliederung (Schülerfrage, -Fälle, -Form, -Weg, -Skizze; oder Frage) – Einstieg nur von dort, nie aus der Bank (pruefung.md § 5, 12.10.)
% 2022-bebb-gk-B4c-m1 (Original 2022-bebb-gk-B4c, Schwere I, rang 9; herausgelöste Fassung: keine)
\begin{pfaufg}{GK~’22\\B4c}{9.}{}
In einem Paketzentrum werden täglich mehrere Millionen Pakete umgeschlagen. Von allen Paketen gehen 12 \% zum Ziel A und 6 \% zum Ziel B, die übrigen an andere Ziele. Unter 100 zufällig ausgewählten Paketen haben genau acht das Ziel B. Berechne, um wie viel Prozent diese Anzahl vom Erwartungswert abweicht.
\karo{2}
\fusshilfe{\mbox{9:~33,3 \%}}
\end{pfaufg}
% 2026-bb-gk-A1.9a-m1 (Original 2026-bb-gk-A1.9a, Schwere III, rang 10; herausgelöste Fassung: keine)
\begin{pfaufg}{GK~’26\\A1.9a}{10.}{\pfohnetr}
Die Zufallsgröße $X$ ist binomialverteilt mit dem Erwartungswert $\mu = 36$ und der Standardabweichung $\sigma = 3$. Bestimme einen Term ohne Parameter, mit dem man $P(X = 37)$ berechnen kann.
\karo{3}
\fusshilfe{\mbox{10:~Term (48 über 37) · 0,75$^{37}$ · 0,25$^{11}$; 0,128}}
\end{pfaufg}

\pfblattende{10}{26 & 25 & 24 & 23 & 22 & 21 & 20 & 19 & 18 & 17}{A1.6a A1.6b A1.9a B4c & A1.9a A1.9b B4c &  & A1.7a A1.7b & B4c &  & B4.2d & B4.2c & B3.2b & }%
{\pfausschnitt{Bernoulli-Kette \pkt{} Binomialwahrscheinlichkeit}{Säulendiagramm und Erwartungswert}{}}
\end{document}